\title{xLSTM: Extended Long Short-Term Memory}

\begin{abstract}
The development of the Long Short-Term Memory (LSTM) framework in the 1990s, notably through the introduction of the constant error carousel and gating mechanisms, laid the groundwork for advancements in sequence modeling. Over subsequent decades, LSTMs have consistently demonstrated effectiveness, serving as key components in a variety of deep learning breakthroughs, and providing the architecture for some of the earliest Large Language Models (LLMs). The rise of Transformers, driven by their parallelizable self-attention mechanism, heralded a shift in dominant paradigms, surpassing LSTMs as the architectures of choice for scaling up large models. In this context, we pose the following question: What level of performance can be achieved in language modeling when LSTMs are scaled to billions of parameters, augmented with state-of-the-art LLM techniques, while systematically addressing traditional shortcomings of LSTMs? To address this, we present exponential gating, accompanied by suitable normalization and stabilization methods. Furthermore, we refine the LSTM’s memory system, yielding: (i) sLSTM, featuring scalar memory, a scalar update mechanism, and an improved approach to memory mixing, and (ii) mLSTM, which introduces a fully parallelizable architecture with matrix-based memory and a covariance update scheme. When these enhanced LSTM modules are deployed within residual block frameworks, we obtain xLSTM blocks, which can be hierarchically stacked to construct xLSTM architectures. The incorporation of exponential gating and innovative memory architectures endows xLSTMs with the capacity to compete strongly with current state-of-the-art Transformer models and State Space Models in both accuracy and scalability.\\
Code available at: \url{https://github.com/NX-AI/xlstm}
\end{abstract}

\section{Introduction}
\label{sec:intro}

The Long Short-Term Memory (LSTM) ideas~\citep{Hochreiter:91,Hochreiter:97nips,Hochreiter:97}, i.e., the constant error carousel and gating, were introduced to overcome the vanishing gradient problem of recurrent neural networks~\citep{Hochreiter:91,Hochreiter:00book}:

\begin{align}
\label{eq:lstm_idea}
  c_t \ = \ f_t \ c_{t-1} \ + \ 
          i_t \  z_t \ , \quad  
          h_t \ = \ o_t \ \psi({c_t}) \ . 
\end{align}

The constant error carousel involves updating the cell state $c_{t-1}$ (indicated in green) through the addition of the cell inputs $z_t$, with this process regulated by sigmoid gating mechanisms (blue). The input gate $i_t$ and forget gate $f_t$ jointly determine how this update proceeds, whereas the output gate $o_t$ regulates the information emitted by the memory cell, specifically, the hidden state $h_t$. Prior to being gated for output, the cell state undergoes normalization or a squashing function $\psi$, after which the resulting value forms the hidden state.

LSTMs have demonstrated substantial success across a diverse array of fields \citep{Hochreiter:01,Hochreiter:07,Schmidhuber:15}, maintaining dominance in text generation applications prior to the emergence of Transformers in 2017~\citep{Vaswani:17}. Their proficiency has been established in a wide range of sequence-related problems, including but not limited to text generation~\citep{Graves:13,Karpathy:15blog}, handwriting synthesis~\citep{Graves:13}, sequence-to-sequence translation~\citep{Sutskever:14nips}, program evaluation~\citep{Zaremba:14arxiv}, creation of image captions~\citep{Karpathy:15,Hossain:19}, code synthesis~\citep{Karpathy:15blog}, modeling rainfall-runoff dynamics~\citep{Kratzert:18,Kratzert:19}, as well as hydrological prediction for flood warnings~\citep{Nearing:24}. In the reinforcement learning landscape, LSTMs have consistently been among the highest-performing sequence models, serving as foundational components in notable systems such as the AlphaStar agent for StarCraft II~\citep{Vinyals:17short}, OpenAI Five for Dota 2~\citep{OpenAI:19}, and magnetic controller models for nuclear fusion experiments~\citep{Degrave:22short}. One of the core strengths of LSTMs is their capacity for abstraction, particularly the extraction and storage of semantic information in their internal states~\citep{Karpathy:15blog}. This capability is illustrated by findings such as number and syntax specialized neurons~\citep{Lakretz:19}, neurons involved in linguistic processing~\citep{Bau:19}, and sentiment-tracking neurons~\citep{Radford:17arxiv}. LSTMs continue to be employed in many significant contemporary domains~\citep{Degrave:22short,Nearing:24}, and have demonstrated enduring relevance over time.

Although LSTMs have achieved remarkable advancements, they are subject to three principal drawbacks: (i) Inability to modify previously made storage choices. This issue is illustrated using the {\it Nearest Neighbor Search} task (see also Appendix~\ref{sec:appNearestNeighborSearch}): Given a reference vector, the model must iteratively examine a sequence to identify the closest matching vector and subsequently return its associated value at the sequence’s conclusion. The left side of Figure~\ref{fig:lstmProblems} reports the mean squared error for this task. LSTMs find it challenging to overwrite a previously stored value when a new, closer vector is encountered, whereas the proposed xLSTM addresses this by using exponential gating mechanisms. (ii) Restricted memory capacity—LSTMs are forced to compress information into scalar-valued cell states. This constraint is highlighted via the {\it Rare Token Prediction} task. The right side of Figure~\ref{fig:lstmProblems} presents perplexity scores from token prediction experiments on Wikitext-103~\citep{Merity:17}, segmented by token frequency. LSTM performance on infrequent tokens suffers as a result of these storage constraints, a challenge our xLSTM overcomes by leveraging a matrix-based memory structure. (iii) Inability to process input in parallel caused by the interdependence of hidden states across time steps—i.e., hidden-to-hidden connections—necessitating strictly sequential computation. 

The deficiencies in LSTM design have contributed to the ascendancy of Transformers~\citep{Vaswani:17} within the realm of language modeling. Given solutions to these shortcomings and scaling LSTMs up to the scale of contemporary Large Language Models, it is natural to ask: how effective could LSTMs become for language modeling under these advancements?

\section{Extended Long Short-Term Memory}
\label{sec:xLSTM}

To address the shortcomings of standard LSTM, Extended Long Short-Term Memory (xLSTM) builds upon the concept in Equation~\eqref{eq:lstm_idea} through two significant enhancements: exponential gating and new memory mechanisms. These developments expand the LSTM family with two additional variants: (i) the sLSTM (see Section~\ref{sec:sLSTM}), which utilizes scalar memory, employs scalar updates, and incorporates a mechanism for memory mixing; and (ii) the mLSTM (see Section~\ref{sec:mLSTM}), which is based on matrix memory with updates following a covariance (outer product) rule that permits full parallelization. The sLSTM and mLSTM models both incorporate exponential gating to strengthen the LSTM framework. To facilitate efficient parallel computation, the mLSTM eschews memory mixing—in other words, it does not utilize hidden-to-hidden recurrent connections. Both sLSTM and mLSTM support extensions to configurations with multiple memory cells; in such cases, sLSTM retains memory mixing across these cells. Moreover, sLSTM can be configured with multiple heads, whereby memory mixing occurs within each head across its cells, but not between the heads. The addition of multiple heads to sLSTM, together with the exponential gating mechanism, enables a novel form of memory mixing. In the case of mLSTM, the distinctions between multiple heads and multiple cells are functionally the same.

By embedding these new LSTM types within residual block modules, xLSTM blocks are formed (refer to Section~\ref{sec:xLSTMblock}). By stacking xLSTM blocks using residual connections, one constructs full xLSTM architectures (see Section~\ref{sec:xLSTMarchitecure}). An overview of the xLSTM architecture and its building blocks is depicted in Figure~\ref{fig:xlstm_sketch}.

\subsection{Review of the Long Short-Term Memory}
\label{sec:orgLSTM}

The initial formulation of the LSTM architecture~\citep{Hochreiter:91,Hochreiter:97nips,Hochreiter:97} established the scalar memory cell as its key storage and computational unit, addressing the issue of vanishing gradients~\citep{Hochreiter:91,Hochreiter:00book} via the constant error carousel mechanism in the cell state update. This memory cell features three gate mechanisms—input, output, and forget gates—with the forget gate introduced later by \citet{Gers:00}. The procedure for updating the LSTM memory cell at a specific time step $t$ is given by:

\begin{align}
c_t \ &= \  f_t \ c_{t-1} \ + \ i_t \ z_t &  & &\text{cell state} \\
h_t  \ &= \ o_t \ \tilde{h}_t \ , & \tilde{h}_t \ &= \ \psi \left( c_t\right)
  &\text{hidden state} \\
z_t \ &= \ \varphi \left( \tilde{z}_t \right) \ , 
  &\tilde{z}_t \ &=  \ \mathbf{w}^\top_{z} \ \mathbf{x}_t \ + \
  r_{z}  h_{t-1} \ + \  b_{z} \ \
  &\text{cell input} \\
i_t \ &= \ \sigma \left( \tilde{i}_t  \right) \ , 
  &\tilde{i}_t \ &= \ \mathbf{w}^\top_{i} \ \mathbf{x}_t \ + \
  r_{i}  \ h_{t-1} \ + \  b_{i} \ \
  &\text{input gate} \\
f_t \ &= \ \sigma \left(  \tilde{f}_t \right) \ , 
  &\tilde{f}_t \ &= \ \mathbf{w}^\top_{f} \ \mathbf{x}_t  \ + \
  r_{f}  \ h_{t-1} \ + \  b_{f} \ \
  &\text{forget gate} \\
o_t \ &= \ \sigma \left( \tilde{o}_t \right) \ , 
  &\tilde{o}_t  \ &= \ \mathbf{w}^\top_{o} \ \mathbf{x}_t \ + \
  r_{o}  \ h_{t-1} \ + \  b_{o} \ \
  &\text{output gate} 
\end{align}

The weight vectors $\mathbf{w}_{z}$, $\mathbf{w}_{i}$, $\mathbf{w}_{f}$, and $\mathbf{w}_{o}$ represent the connections from the input vector $\mathbf{x}_t$ to the cell input, input gate, forget gate, and output gate, respectively. The parameters $r_{z}$, $r_{i}$, $r_{f}$, and $r_{o}$ denote the recurrent weights linking the previous hidden state $h_{t-1}$ with the corresponding cell input, input gate, forget gate, and output gate. The bias terms, $b_{z}$, $b_{i}$, $b_{f}$, and $b_{o}$, are associated with these operations. The cell input and hidden state activation functions are denoted by $\varphi$ and $\psi$—the latter typically being $\tanh$—where $\psi$ serves to bound the cell state that would otherwise be unregulated. The activation functions applied at each gate are sigmoid, explicitly, $\sigma(x) = 1/(1 + \exp(-x))$. In subsequent versions, the architecture was extended to support vectors of memory cells $\mathbf{c}_t \in \mathbb{R}^d$, allowing the implementation of recurrent weight matrices $\mathbf{R} \in \mathbb{R}^{d \times d}$ to blend the outputs from different memory cells~\citep{Greff:15}; further elaboration is available in Appendix~\ref{sec:apporgLSTM}. Investigations using ablation analysis have demonstrated the necessity of each memory cell component~\citep{Greff:15}.

\subsection{sLSTM}
\label{sec:sLSTM}

In order to enhance LSTMs with the capacity to adjust their storage choices, we incorporate exponential gates (illustrated in red), in combination with mechanisms for normalization and stabilization. Specifically, the input and forget gates are allowed to use exponential activation functions. For the normalization process, a normalizer state is proposed, which accumulates the sum of the products between the input gate and each successive forget gate.

The scalar sLSTM forward pass is:

\begin{align}
\label{eq:slstmforward}
c_t \ &= \  f_t \ c_{t-1} \ + \ i_t \ z_t & & 
 &\text{cell state} \\
n_t \ &= \  f_t \ n_{t-1} \ + \ i_t  & & 
 &\text{normalizer state} \\
h_t \ &= \ o_t \ \tilde{h}_t \ ,
  &\tilde{h}_t \ &= \ c_t / n_t
&\text{hidden state} \\
z_t \ &= \ \varphi \left( \tilde{z}_t \right) \ , 
  &\tilde{z}_t \ &=  \ \mathbf{w}^\top_{z} \ \mathbf{x}_t \ + \
  r_{z}  h_{t-1} \ + \  b_{z}
  &\text{cell input} \\
i_t \ &= \ \!\exp \left( \tilde{i}_t  \right) \ , 
  &\tilde{i}_t \ &= \ \mathbf{w}^\top_{i} \ \mathbf{x}_t \ + \
  r_{i}  \ h_{t-1} \ + \  b_{i}
  &\text{input gate} \\
f_t \ &= \ \sigma \left(  \tilde{f}_t \right) \ \text{OR} \
\!\exp \left(  \tilde{f}_t \right) \ ,
  &\tilde{f}_t \ &= \ \mathbf{w}^\top_{f} \ \mathbf{x}_t  \ + \
  r_{f}  \ h_{t-1} \ + \  b_{f}
  &\text{forget gate} \\
o_t \ &= \ \sigma \left( \tilde{o}_t \right) \ , 
  &\tilde{o}_t  \ &= \ \mathbf{w}^\top_{o} \ \mathbf{x}_t \ + \
  r_{o}  \ h_{t-1} \ + \  b_{o}
  &\text{output gate} 
\end{align}

We transfer the original LSTM gating techniques, i.e., input- and/or hidden-dependent gating plus bias term, to the new architectures. Exponential activation functions can lead to large values that cause overflows. Therefore, we stabilize gates with an additional state \mbox{$m_t$~\citep{Milakov:18arxiv}}:

\begin{align}
\label{eq:slstmstabil}
m_t \ &= \ \max \left( \log ( f_t ) + m_{t-1} , \log ( i_t ) \right) &\text{stabilizer state} \\
i'_t \ &= \ \!\exp \left( \log \left ( i_t \right) - m_t \right) \ = \ \!\exp \left( \tilde{i}_t - m_t \right) \ \, 
  &\text{stabil. input gate} \\
f'_t \ &= \ \!\exp \left( \log \left( f_t \right) + m_{t-1} - m_t \right) \ \, 
  &\text{stabil. forget gate}
\end{align}

Appendix~\ref{sec:appsLSTM} demonstrates that substituting $f_t$ with $f'_t$ and $i_t$ with $i'_t$ during the forward computation has no effect on either the network’s output or the gradients of the loss with respect to the parameters.

\paragraph{New Memory Mixing.} Like the standard LSTM, sLSTM is capable of utilizing several memory cells (see Appendix~\ref{sec:appsLSTM}). The presence of multiple cells facilitates memory mixing through the use of recurrent connections, namely $\mathbf{R}_{\mathbf{z}}$, $\mathbf{R}_{\mathbf{i}}$, $\mathbf{R}_{\mathbf{f}}$, and $\mathbf{R}_{\mathbf{o}}$, which connect the hidden state vector $\mathbf{h}$ with the memory cell input $\mathbf{z}$, as well as with the input, forget, and output gates $\mathbf{i}$, $\mathbf{f}$, and $\mathbf{o}$. An important innovation in this context arises from the utilization of exponential gating in memory mixing. Within the updated sLSTM framework, multiple heads can be assigned, permitting memory mixing among memory cells within each individual head while prohibiting inter-head mixing. The combination of sLSTM heads and exponential gating thus introduces a novel mechanism for memory interaction.

\subsection{mLSTM}
\label{sec:mLSTM}

To improve the storage capacity of LSTMs, we substitute the scalar memory cell $c \in \mathbb{R}$ with a memory matrix $\mathbf{C} \in \mathbb{R}^{d \times d}$. As a result, information retrieval involves multiplying by this matrix. At each time step $t$, our goal is to store a key vector $\mathbf{k}_t \in \mathbb{R}^d$ along with an associated value vector $\mathbf{v}_t \in \mathbb{R}^d$, adopting notation used in the Transformer architecture. Subsequently, at a later time $t + \tau$, retrieval of $\mathbf{v}_t$ is achieved using a query vector $\mathbf{q}_{t+\tau} \in \mathbb{R}^d$. This framework aligns with the model of Bidirectional Associative Memories (BAMs) \citep{Kohonen:72,Anderson:72,Nakano:72,Anderson:77}.

The covariance update rule~\citep{Sejnowski:77,Dayan:91} for storing a key-value pair is

\begin{align}
\mathbf{C}_t \ &= \ \mathbf{C}_{t-1} \ + \ \mathbf{v}_{t} \ \mathbf{k}_t^\top \ .
\end{align}

We posit that applying a layer-norm prior to the projection of input vectors into key and value spaces ensures these vectors exhibit zero mean. The covariance-based update mechanism achieves optimality~\citep{Dayan:91} in maximizing the separability of the retrieved binary representations—this optimum aligns with achieving a maximal signal-to-noise ratio. More pronounced separability can be realized if retrieval is constrained to pairwise interactions, albeit at the cost of quadratic computational complexity as encountered in attention mechanisms~\citep{Krotov:16,Krotov:17,Ramsauer:21}. Notably, the covariance update approach aligns with the methodology in Fast Weight Programmers~\citep{Schmidhuber:92ncfastweights,Schlag:21}, which were subsequently extended by introducing a fixed decay term applied to $\mathbf{C}_{t-1}$ and a fixed learning rate for scaling 
$\mathbf{v}_{t} \mathbf{k}_t ^\top$~\citep{Ba:16}. Drawing from this perspective, we adapt the covariance update strategy within the LSTM architecture: here, the forget gate modulates decay, the input gate governs the learning rate, and the output gate is responsible for the scaling of the retrieved vector.

For this matrix memory, the normalizer state is the weighted sum of key vectors, where each key vector is weighted by the input gate and all future forget gates. Again, the normalizer state keeps record of the strength of the gates. Since the dot product between query and normalizer state can be close to zero, we use the absolute value of this dot product and lower bound it by a threshold (typically 1.0) as done previously~\citep{Sun:23arxiv}. The mLSTM forward pass is:

\begin{align}
\label{eq:mlstm_recurrent_begin}
\mathbf{C}_t \ &= \  f_t \ \mathbf{C}_{t-1} \ + \ 
  i_t \ \mathbf{v}_t \ \mathbf{k}_t^\top & &  &\text{cell state} \\
\mathbf{n}_t \ &= \  f_t \ \mathbf{n}_{t-1} \ + \ 
  i_t \ \mathbf{k}_t & &  &\text{normalizer state} \\
\mathbf{h}_t  \ &= \ \mathbf{o}_t \ \odot \ \tilde{\mathbf{h}}_t
\ , \qquad \qquad 
\tilde{\mathbf{h}}_t \ = \   \mathbf{C}_t \mathbf{q}_t \ / \ 
  \max \left\{ |\mathbf{n}_t^\top \mathbf{q}_t|, 1 \right\} 
   &&&\text{hidden state} \\
\mathbf{q}_t \ &= \ \mathbf{W}_q \ \mathbf{x}_t \ + \ \mathbf{b}_q  & &  &\text{query input} \\
\mathbf{k}_t \ &= \ \frac{1}{\sqrt{d}} \mathbf{W}_k \ \mathbf{x}_t \ + \ \mathbf{b}_k  & &  &\text{key input} \\
\mathbf{v}_t \ &= \ \mathbf{W}_v \ \mathbf{x}_t \ + \ \mathbf{b}_v  & &  &\text{value input} \\
i_t \ &= \ \!\exp \!  \! \left( \tilde{i}_t  \right) \ , 
 \qquad \qquad \ \ \ \ \,
  \tilde{i}_t \ = \ \mathbf{w}^\top_{i} \ \mathbf{x}_t \ + \  b_{i} \ \ \ 
  &&&\text{input gate} \\
f_t \ &= \ \sigma \!  \left(  \tilde{f}_t \right) \ \text{OR} \ \!\exp \!  \! \left(  \tilde{f}_t \right) , \ \ \: \!
\tilde{f}_t \ = \ \mathbf{w}^\top_{f} \ \mathbf{x}_t  \ + \
  b_{f}
  &&& \text{forget gate} \\
\label{eq:mlstm_recurrent_end}
\mathbf{o}_t \ &= \ \sigma \left( \tilde{\mathbf{o}}_t \right) \ , \qquad \qquad \quad \ \ \ \,
  \tilde{\mathbf{o}}_t  \ = \ \mathbf{W}_{\mathbf{o}} \ \mathbf{x}_t \ + \
  \mathbf{b}_{\mathbf{o}}
  &&&\text{output gate} 
\end{align}

The mLSTM architecture supports several memory cells in the same manner as the conventional LSTM. In mLSTM, there is an equivalence between deploying multiple heads and multiple cells, due to the absence of any interaction or mixing between memory components. To address potential instability from exponential gating mechanisms in mLSTM, we employ the identical stabilization strategies utilized for sLSTM (refer to Equation~\ref{eq:slstmstabil}). 
Given that mLSTM does not mix its memory states, its recurrence relation can be represented in a parallelized format. Further information can be found in Appendix~\ref{sec:appmLSTM}.

\subsection{xLSTM Architecture}
\label{sec:xLSTMarch}

\paragraph{xLSTM Blocks.}
\label{sec:xLSTMblock}
The function of an xLSTM block is to generate a nonlinear summary of the historical sequence data within a high-dimensional representation, thereby improving the ability to distinguish among various sequence histories or contexts. This enhanced separation of historical patterns is essential for accurately forecasting subsequent elements in the sequence, such as the next token. Our approach is informed by Cover’s Theorem~\citep{Cover:65}, which posits that patterns that have undergone nonlinear embedding in a high-dimensional space are more amenable to linear separation than those in the lower-dimensional original space. We examine two main designs for the residual block: (i) A residual block that employs post up-projection (analogous to the approach used in Transformers), which first applies a nonlinear summarization within the initial feature space, subsequently projects the result linearly into a higher-dimensional space, introduces a non-linear activation, and then reduces the dimensionality back to the original space via a linear transformation; this process is illustrated in the leftmost panel of Figure~\ref{fig:Backbones} and in the third column of Figure~\ref{fig:xlstm_sketch}. For greater clarity, a comprehensive version is available in Figure~\ref{fig:appsLSTM_detailed} in the appendix. (ii) Alternatively, we consider a residual block with pre up-projection (reminiscent of State Space Models), which begins by performing a linear projection into a higher-dimensional space,

performs a non-linear aggregation of historical information in a high-dimensional feature space, subsequently employing a linear transformation to project the representation back to the input space. For xLSTM blocks that incorporate an sLSTM, the architecture includes a post up-projection module. Conversely, for those featuring an mLSTM, a pre up-projection block is adopted due to the increased memory footprint in the elevated dimensional space. Additional illustrations and explanation are provided in the leftmost panel of Figure~\ref{fig:Backbones}, the third column of Figure~\ref{fig:xlstm_sketch}, and Figure~\ref{fig:appsLSTM_detailed} within the appendix.

\paragraph{xLSTM Architecture. }
\label{sec:xLSTMarchitecure}
The xLSTM architecture is formed by stacking its constituent blocks in a residual fashion, as described in prior work~\citep{Srivastava:15,He:16}. This design leverages pre-LayerNorm residual backbones~\citep{Ba:16b}, consistent with practices prevalent in modern Large Language Models. Refer to the final column of Figure~\ref{fig:xlstm_sketch} for an architectural overview.

\subsection{Memory and Speed Considerations}

In contrast to Transformers, xLSTM architectures maintain both linear computational cost and fixed memory requirements as the sequence length increases. Owing to their compressive memory design, xLSTM models are highly advantageous for use in industrial scenarios and in edge device deployments.

mLSTM networks employ a parameter-free memory structure, yet the associated computations are intensive, owing to the need for $d \times d$ matrix storage and corresponding $d \times d$ updates. Our approach navigates the balance between expanding memory capability and minimizing computational load. Nevertheless, since these calculations are amenable to GPU parallelization, their impact on actual execution time remains relatively minor.

Whereas mLSTM operations support parallelization similar to FlashAttention~\citep{Dao:22,Dao:23} and GLA~\citep{Yang:23arxiv}, sLSTM models face constraints due to memory mixing, arising from hidden-to-hidden interactions, which hinder parallel processing. To address this, we introduced a specialized CUDA-based implementation optimized for register-level GPU memory use, resulting in performance that is generally less than twice as slow as mLSTM.

\section{Related Work}
\label{sec:related}

\paragraph{Linear Attention.} A variety of strategies have been developed to address the Transformer's quadratic computational demands relative to context length, aiming to achieve linear scaling instead. For instance, the Synthesizer dispenses with explicit token-to-token interactions by directly learning synthetic attention weights~\citep{Tay:20arxiv}. Linformer approaches self-attention through a low-rank matrix factorization, enabling a linear approximation of the mechanism~\citep{Wang:20arxiv}. Linear Transformer modifies the attention process to admit a linear computation in practice \citep{Katharopoulos:20}. Performer achieves an efficient linear approximation of softmax attention via the use of positive orthogonal random features~\citep{Choromanski:21}. Furthermore, attention modules have been substituted by rapid, long-range convolution operations in models such as the Structured Global Convolution (SGConv)~\citep{Li:22} and the Hyena Hierarchy~\citep{Poli:23}.

\paragraph{State Space Models.} In recent years, State Space Models (SSMs) have attracted significant attention due to their linear computational complexity relative to input sequence length and their competitive results against Transformer architectures. Early contributions to this area include the Structured State Space sequence model (S4)~\citep{Gu:21}. Subsequent advancements were seen in models such as the Diagonal State Space (DSS) model~\citep{Gupta:22}, Gated State Space (GSS) models~\citep{Mehta:22}, the S5 model~\citep{Smith:22}, Bidirectional Gated SSM (BiGS)~\citep{Wang:22}, the H3 model~\citep{Fu:23}, and, more recently, Mamba~\citep{Gu:24arxiv}.

\paragraph{Recurrent Neural Networks.} Due to their ability to process sequences with linear time complexity relative to context length, Recurrent Neural Networks (RNNs) have been explored as alternatives to the Transformer and attention mechanisms. Notably, RNNs equipped with Deep Linear Recurrent Units (LRUs) have achieved encouraging outcomes on language modeling tasks~\citep{Orvieto:23, De:24arxiv}. Similarly, the Hierarchically Gated Linear RNN (HGRN)~\citep{Qin:23} and its successor HGRN2~\citep{Qin:24arxiv} have demonstrated efficacy in this area. Among prominent RNN-based models for large-scale language modeling, RWKV~\citep{Peng:23arxivshort,Peng:24arxivshort} stands out for delivering performance on par with Transformer-based approaches.

\paragraph{Gating.}
Gating, initially developed within the LSTM framework, has since become a central component in numerous contemporary models, often revisited and adapted in different contexts. Architectures employing gating include HGRN~\citep{Qin:23}, HGRN2~\citep{Qin:24arxiv}, Gated Linear Attention (GLA)~\citep{Yang:23arxiv}, Gated State Space (GSS) models~\citep{Mehta:22}, Bidirectional Gated SSM (BiGS)~\citep{Wang:22}, Moving Average Equipped Gated Attention (MEGA)~\citep{Ma:22}, RWKV~\citep{Peng:23arxivshort}, and Mamba~\citep{Gu:24arxiv}.

\paragraph{Covariance Update Rule.} 
In order to improve storage efficiency, our mLSTM cell utilizes a matrix-based memory governed by a covariance update procedure. Related strategies that apply similar update principles can be found in Fast Weight Programmers~\citep{Schmidhuber:92ncfastweights,Schlag:21}, RWKV-5 and RWKV-6~\citep{Peng:24arxivshort}, Retention~\citep{Sun:23arxiv}, Linear Transformer~\citep{Katharopoulos:20}, and HGRN2~\citep{Qin:24arxiv}.

\paragraph{Most Related.} From a conceptual standpoint, the models most similar to xLSTM include Retention~\citep{Sun:23arxiv}, RWKV~\citep{Peng:23arxivshort,Peng:24arxivshort}, and HGRN2~\citep{Qin:24arxiv}, all of which incorporate either a matrix-based memory structure or some form of gating mechanism. Distinct from the recently introduced sLSTM, these prior methods lack support for memory mixing. The ability to mix memories is critical for addressing state tracking challenges, consequently rendering LSTMs more expressive than both State Space Models (SSMs) and Transformers~\citep{Merrill:24,Deletang:23}. State tracking is essential in contexts such as code interpretation or entity tracking over extended narratives.

\paragraph{Residually Stacking Architectures.} Consistent with the design of most state-of-the-art deep learning systems, xLSTM architectures employ a strategy of residual stacking of their fundamental components~\citep{Srivastava:15,He:16}.

This architectural paradigm has proven crucial for the success of deep convolutional networks~\citep{He:16} and, more notably, for Transformers~\citep{Vaswani:17}. Transformers serve as the foundational architecture for current Large Language Models (LLMs) such as GPT-3~\citep{Brown:20short}, ChatGPT~\citep{Schulman:22short}, GPT-4~\citep{Achiam:23gpt4short}, Megatron-LM~\citep{Shoeybi:19}, Gopher~\citep{Rae:21short}, ERNIE 3.0 Titan~\citep{Wang:21arxivshort}, GLaM ~\citep{Du:21glamshort}, Chinese M6~\citep{Lin:21}, multilingual AlexaTM 20B~\citep{Soltan:22}, OPT~\citep{Zhang:22opt}, Chinchilla~\citep{Hoffmann:22short}, BLOOM~\citep{Scao:22arxivshort}, GLM-130B~\citep{Zeng:22arxivshort}, LaMDA~\citep{Thoppilan:22short}, PaLM~\citep{Chowdhery:22short}, Llama~\citep{Touvron:23arxiv}, and Gemini~\citep{GeminiTeam:23arxiv,Reid:24arxiv}.

\section{Experiments}
\label{sec:Exp}

This section presents an empirical assessment of xLSTM, positioning its language modeling performance relative to established techniques. Section~\ref{sec:ExpSyntethic} centers on probing the unique strengths of xLSTM by applying it to purposefully constructed synthetic problems. In Section~\ref{sec:ExpComparison}, we report the validation perplexity of several modern language modeling approaches, all trained on 15B SlimPajama~\citep{Soboleva:23} tokens. Within this context, we also include an ablation study on xLSTM. We subsequently examine the scaling properties of each approach following methodologies in \citet{Kaplan:20} and \citet{Brown:20short}. 

Section~\ref{sec:ExpLanguage} expands the analysis to a more extensive language modeling setting. Here, both xLSTM and the top competitors from Section~\ref{sec:ExpComparison} are trained with an order of magnitude larger data volume: 300B SlimPajama~\citep{Soboleva:23} tokens. We first gauge the capacity of these models to generalize over extended contexts, then assess performance through validation set perplexity and downstream task outcomes~\citep{Sutawika:24}. Next, we measure effectiveness across the 571 text domains found in the PALOMA benchmark dataset~\citep{Magnusson:23arxivshort}. Finally, we revisit the scaling analysis for each method, now under conditions involving twentyfold more training examples.

\label{sec:xlstm_blocks_notation}
Throughout our experiments, we denote xLSTM[$a$:$b$] as representing a configuration in which the proportion of mLSTM-based to sLSTM-based xLSTM blocks is $a/b$. For instance, xLSTM[7:1] specifies that, among eight total blocks, seven employ mLSTM and one utilizes sLSTM. In scenarios involving 48 blocks in total, this composition results in 42 blocks based on mLSTM and 6 based on sLSTM. Additionally, in all experiment setups, we incorporate up-projection blocks before mLSTM modules and after sLSTM modules, respectively.

\subsection{Synthetic Tasks and Long Range Arena}
\label{sec:ExpSyntethic}
To begin, we evaluate xLSTM’s novel exponential gating combined with memory mixing using formal language benchmarks~\citep{Deletang:23}. Next, we investigate how xLSTM’s updated matrix memory performs on the Multi-Query Associative Recall benchmark~\citep{Arora:23arxiv}. Lastly, the capability of xLSTM for handling lengthy sequences is measured using the Long Range Arena~\citep{Tay:21}.

\paragraph{Evaluating xLSTM's Exponential Gating Incorporating Memory Mixing.} We assess the performance of xLSTM's novel exponential gating mechanism with memory mixing, which is designed to address state-tracking challenges~\citep{Merrill:24, Merrill:23}. Our evaluation builds on and enhances the formal language benchmarks from \citet{Deletang:23}, allowing for multi-length training that facilitates extrapolation to sequences of greater lengths. For comprehensive task descriptions and more exhaustive experimental results, consult Appendix~\ref{sec:appExpSynthetic-formalLang}. 

We conduct a comparative study between xLSTM and alternative architectures, encompassing Transformers, State Space Models, and classic Recurrent Neural Networks. To measure performance, we compute accuracy over task-relevant tokens, normalizing this metric such that 0 indicates chance-level and 1 indicates perfect prediction. Our experiments utilize 2-block versions of various architectures, including xLSTM[0:1] (representing sLSTM only), xLSTM[1:0] (representing mLSTM only), xLSTM[1:1], Llama, Mamba, RWKV, Retention, Hyena, standard LSTM, and LSTM integrated within Transformer blocks (LSTM (Block)). Figure~\ref{fig:formal-main} displays the empirical findings of these comparisons.

Importantly, models lacking memory mixing—such as Transformers and State Space Models that do not perform state tracking—are unable to solve tasks involving, for instance, regular grammars like the parity task. These outcomes are consistent with previous reports highlighting the inherent limitations of Transformers and State Space Models relative to RNNs~\citep{Merrill:24, Merrill:23, Deletang:23}.

\paragraph{Test of xLSTM's Memory Capacities on Associative Recall Tasks.} This study examines the matrix-based memory mechanisms introduced in xLSTM, focusing on its ability to store and retrieve information in the Multi-Query Associative Recall task~\citep{Arora:23arxiv}. In this setup, each input sequence is composed of key-value pairs that are sampled at random from an extensive vocabulary, requiring the model to retain these associations for subsequent access. To intensify the challenge posed by the benchmark, we have scaled up both the number of key-value pairs—reaching as many as 256—and the overall context length, which extends to 2048 tokens. This adjustment provides a more rigorous evaluation of the memory limits across various architectures. Our comparison includes two-block variants of Llama, Mamba, RWKV-5, RWKV-6, xLSTM[1:1], and xLSTM[1:0]. We assess these models based on the accuracy of recalling the relevant pairs. As the memory size in Transformers such as Llama scales exponentially with the dimensionality of the representations~\citep{Ramsauer:21}, such models set a benchmark standard for this task. The empirical findings can be found in Figure~\ref{fig:mqar-main}. Notably, among non-Transformer approaches, xLSTM[1:1] attains the highest performance, even when model size is restricted. Contrary to possible expectations, the inclusion of an sLSTM block not only avoids reducing memory efficacy, but instead augments it, with the improvement especially prominent under the scenario involving 256 key-value associations. Further experimental outcomes are provided in Appendix~\ref{sec:appExpSynthetic-mqar}, where extrapolation studies suggest that the advanced memory structure of xLSTM facilitates generalization to much longer contexts than those present in the training set.

\paragraph{Test of xLSTM's Long Context Capabilities on Long Range Arena.} To evaluate xLSTM's effectiveness in processing extended sequences and sizable contexts, we benchmark various techniques on the Long Range Arena~\citep{Tay:21}. xLSTM consistently achieves strong results across every task, indicating that its architecture adeptly manages the diverse challenges presented by long context scenarios.

For more details, see Appendix~\ref{sec:appExpSynthetic-lra}.

\subsection{Method Comparison and Ablation Study}
\label{sec:ExpComparison}

In order to explore the central question of this study—specifically, the capabilities of our newly proposed LSTM variants at scale for language modelling—we conduct experiments in which xLSTMs, Transformers, State Space Models, and various other techniques are trained using 15B tokens sampled from SlimPajama, employing an identical auto-regressive training framework. Model performance is then evaluated on a shared validation set, and we carry out ablation studies to further analyze the behavior of xLSTMs.

\paragraph{Comparison of xLSTM with Alternative Approaches.} To benchmark performance, we trained several models using 15B tokens sourced from the SlimPajama dataset~\citep{Soboleva:23}. Each model's performance was assessed via perplexity on a held-out validation set. The set of models assessed includes our proposed xLSTM, GPT-3 (a Transformer architecture)~\citep{Brown:20short}, Llama (Transformer)~\citep{Touvron:23arxiv}, H3 (a State Space Model, SSM)~\citep{Fu:23}, Mamba (SSM)~\citep{Gu:24arxiv}, RWKV-4, RWKV-5, and RWKV-6 (all Recurrent Neural Networks, RNNs)~\citep{Peng:23arxivshort, Peng:24arxivshort}, GLA (a linear Transformer model)~\citep{Yang:23arxiv}, HGRN and HGRN2 (RNNs)~\citep{Qin:23, Qin:24arxiv}, as well as RetNet and Hyena (linear Transformers)~\citep{Sun:23arxiv, Poli:23}. Variants xLSTM[1:0] and xLSTM[7:1] (notation explained in Section~\ref{sec:xlstm_blocks_notation}) were also included. 

All models were trained using mixed precision. However, RWKV-5, RWKV-6, GLA, and HGRN2 did not rely on the default PyTorch automatic mixed precision approach, as detailed in Appendix Section~\ref{sec:appGenTrainProc}. The methods considered fall into three groups: (a) traditional Transformers, (b) State Space Models (SSMs), and (c) Recurrent Neural Networks (RNNs) alongside linear Transformer architectures, which employ linear approximations in place of standard Transformer attention mechanisms. All experimental models were scaled to match the size of GPT-3 with 350M parameters, featuring an embedding size of 1024 and 24 residual blocks. Only GPT-3 utilizes tied token and output embedding weights, resulting in a modestly lower parameter count.

Table~\ref{tab:spaj15b_model_results} highlights that xLSTM achieves the lowest validation perplexity among all tested models. For further methodological specifics and additional results, the reader is directed to Appendix~\ref{sec:appExpComparison}. The scaling trends for these experiments are presented in Figure~\ref{fig:temp_sclaw15B}, which demonstrate that xLSTM’s relative performance advantage is likely to persist as model capacity increases.

\paragraph{Ablation Studies.}
As illustrated in Table~\ref{tab:spaj15b_model_results} and Figure~\ref{fig:temp_sclaw15B}, xLSTM attains strong language modeling performance when trained on 15B tokens from SlimPajama. To systematically examine the impact of each architectural modification from the standard LSTM to the full xLSTM design, we incrementally transform a basic LSTM architecture into xLSTM. The first modification replaces standard layers with LSTM layers embedded within pre-LayerNorm residual structures. Next, this setup is augmented by including a post up-projection component. The final changes introduce exponential gating and matrix memory to complete the xLSTM architecture.

The results are shown in Table~\ref{tab:ablstudies} (top). 

Ablation experiments indicate that both the exponential gating mechanism and the matrix memory contribute significantly to the observed gains in performance. Given the crucial role of gating within RNNs and State Space Models, we further examine various gating strategies. As presented in Table~\ref{tab:ablstudies} (bottom), our findings demonstrate that enabling each gate to be input-dependent and trainable yields progressively better results. Further exploration of the effects of individual backbone modules is provided in Appendix~\ref{sec:appAblDetails}.

\subsection{xLSTM as Large Language Model}
\label{sec:ExpLanguage}

This work concludes with comprehensive experiments in large-scale language modeling, aiming to evaluate xLSTM's effectiveness as a large language model (LLM). To this end, we significantly scale up the training dataset to 300B tokens drawn from SlimPajama, mirroring the quantity of training data employed by methods such as Mamba~\citep{Gu:24arxiv} and Griffin~\citep{De:24arxiv}.

For our comparative analysis, we pit xLSTM against RWKV-4, Llama, and Mamba—each selected to represent a distinct methodological class as outlined in Section~\ref{sec:ExpComparison}. RWKV-4 is chosen as the benchmark RNN, since optimal training precision configurations for RWKV-5, RWKV-6, and HGRN2 were established only after commencement of the 300B token training experiments (refer to Appendix~\ref{sec:appGenTrainProc}).

Our experimental setup encompasses a range of model scales (125M, 350M, 760M, and 1.3B parameters), and all variants are evaluated for their ability to generalize to longer sequence lengths as well as their validation performance. We further analyze these models' efficacy on downstream tasks, examine their language modeling capabilities across 571 text domains included in the PALOMA benchmark, and systematically probe their scaling law trends.

\paragraph{Sequence Length Extrapolation.}
\label{sec:spaj300B_ctx_extrapolation}
We begin by evaluating the ability of 1.3B parameter versions of xLSTM, RWKV-4, Llama, and Mamba to generalize to longer input sequences. Each model is trained with a context length of 2048 and subsequently assessed on sequences extending up to 16384 tokens. The comparative results are illustrated in Figure~\ref{fig:spaj300B_seq_extrapolation}. Notably, among the evaluated architectures, xLSTM consistently achieves lower perplexity scores over extended contexts.

\paragraph{Validation Perplexity and Downstream Tasks.}
\label{sec:spaj_downstream_eval}
Next, we assess all model variants by measuring their validation perplexity on next token prediction using the SlimPajama benchmark and evaluating their effectiveness on downstream tasks centered on common sense reasoning. These evaluations are conducted across every model scale to compare the relative capabilities of xLSTM, RWKV-4, Llama, and Mamba.

The third column in Table~\ref{tab:lmeval} presents the perplexity values on the validation set for various approaches. For all model sizes, both xLSTM[1:0] and xLSTM[7:1] achieve the lowest validation perplexity, indicating superior performance by these models. The remaining columns in Table~\ref{tab:lmeval} summarize the outcomes on downstream benchmarks. Across nearly all evaluated tasks and model configurations, xLSTM outperforms competing methods; Mamba demonstrates superior results exclusively on certain instances of the ARC task. Further information can be found in Appendix~\ref{sec:appExpLanguage}.

\paragraph{Performance on PALOMA Language Tasks.} As a third evaluation, we compare xLSTM, RWKV-4, Llama, and Mamba across all parameter scales by assessing their next token prediction capabilities on the PALOMA language tasks~\citep{Magnusson:23arxivshort}. The models' effectiveness is quantified by their perplexity scores on next token prediction across 571 diverse text domains, encompassing sources from nytimes.com to Reddit’s r/depression.

Table~\ref{tab:ppleval} presents the token prediction perplexity results, distinguishing between broad language modeling (the initial seven columns) and a selection of granular domain-specific benchmarks (the final five columns). On these language tasks, xLSTM[1:0] consistently achieves lower perplexity compared to xLSTM[7:1].

More specifically, xLSTM[1:0] attains a lower perplexity than Mamba in 568 out of 571 cases (99.5\%), surpasses Llama in 486 out of 571 domains (85.1\%), and outperforms RWKV-4 in 570 out of 571 text domains (99.8\%). Further results are available in Appendix~\ref{sec:appExpLanguage}.

\paragraph{Scaling Laws.} Next, we investigate the presence of power-law scaling, which enables predictions of model performance as size increases~\citep{Kaplan:20,Brown:20short}. Figure~\ref{fig:temp_sclaw300B} illustrates this scaling pattern. While all models demonstrate qualitatively comparable scaling trends, their baselines differ. Among the group, RWKV-4 consistently yields the lowest performance, with Llama and Mamba trailing subsequently. xLSTM outperforms Mamba by approximately the same margin that separates Mamba from Llama. These scaling results suggest that as model size grows, xLSTM is likely to maintain its performance advantage over both Transformer and State-Space architectures.

\textbf{Generation Times and Maximal Throughput.} We evaluate the generation latency in Figure~\ref{fig:inference_time_generation} and report peak decoding throughput in Figure~\ref{fig:inference_time_decoding_throughput} (left) across our xLSTM architectures at the 1.3B parameter level. For comparison, we include equivalent-scale implementations of Mamba, Llama, and RWKV from HuggingFace, employing a static key-value cache for the Llama runs. It should be noted that, at the time these experiments were conducted, neither full cache compilation of the Transformer Model nor applying \texttt{torch.compile} to Mamba was operational. Each generation benchmark uses a batch size of 1 and a pre-fill value of 16 across all model types, a setup particularly advantageous for Transformer-based models.

The results in Figure~\ref{fig:inference_time_generation} highlight that xLSTM, like the other recurrent architectures—Mamba and RWKV-4—exhibits linear computational complexity, in contrast to the quadratic scaling observed in the Llama model. For throughput evaluation, we test a range of batch sizes and pre-fill settings with Llama. As illustrated in Figure~\ref{fig:inference_time_decoding_throughput} (right), xLSTM supports substantially larger batch sizes thanks to its constant memory footprint, thereby achieving superior throughput compared to Llama.

\section{Limitations}

(i) Unlike mLSTM, the way sLSTM mixes memory prevents the use of parallelizable computations, making it unsuited to rapid parallel implementations. However, we implemented an efficient CUDA kernel for sLSTM, which currently runs at less than double the execution time of our parallel mLSTM solution. (ii) The current CUDA kernels for mLSTM lack optimization, which results in a runtime roughly four times slower compared to FlashAttention and the scan operation applied in Mamba. Improvements similar to those used in FlashAttention could lead to faster CUDA kernels. (iii) The requirement to process $d \times d$ matrices in the mLSTM’s matrix memory incurs substantial computational cost. However, both memory updating and reading are parameter-free operations that can leverage standard matrix parallelization techniques, so the increased computation demands have minimal impact on wall clock time. (iv) Careful attention is needed when initializing the forget gates. (v) Because the matrix memory is not a function of sequence length, extending the sequence may saturate the memory capacity when handling much longer contexts. Nevertheless, our findings indicate this is not an issue for sequence contexts up to 16k tokens (refer to Section~\ref{sec:spaj300B_ctx_extrapolation}). (vi) Owing to the computational intensity of running large-scale language experiments, we have not conducted thorough architecture or hyperparameter optimization, especially in the case of larger xLSTM configurations. We expect that comprehensive optimization will be essential for xLSTM to realize its maximum capabilities.

\section{Conclusion}
We have partially addressed our central question: To what extent can language modeling be advanced by scaling LSTM architectures to the order of billions of parameters? At this stage, our findings suggest: ``At least as far as present methodologies such as Transformers or State Space Models''. We have introduced enhancements to the LSTM architecture, resulting in xLSTM, through the application of exponential gating featuring memory mixing and the integration of an updated memory structure. Comparative evaluation demonstrates that xLSTM models achieve competitive, and at times superior, performance in language modeling tasks relative to leading approaches such as Transformers and State Space Models. The identified scaling behavior further implies that larger xLSTM architectures may represent strong alternatives to current Large Language Models that employ Transformer-based designs. Beyond language modeling, the advancements in xLSTM may have significant ramifications in domains including Reinforcement Learning, Time Series Forecasting, and physical systems modeling.

\end{document}